\documentclass[12pt]{article}
\usepackage[utf8]{inputenc}
\usepackage[T1]{fontenc}
\usepackage{hyperref}
\title{Tanibili Language Resources for K-12 Curriculum Development}
\author{Compiled by Claude Sonnet 5 (claude-sonnet-5) for Andr\'as Kornai}
\date{}
\begin{document}
\maketitle

\section{Introduction}
Tanibili (ISO 639-3 tbe; also spelled Tanimbili, and called Nyisunggu or Tanabili) is an Austronesian language of the Utupua-Vanikoro branch of the Temotu group. It is spoken on Utupua in Temotu Province, Solomon Islands. It is not the Santa Cruz (Nat\"ugu) language of Nend\"o. The sources found give about 15 speakers (1999) and a UNESCO classification of severely endangered.

\section{Primary Educational Resources}
No primers, textbooks or school programmes were found.

\section{Linguistic and Reference Materials}
\begin{itemize}
\item Wikipedia article on Tanimbili. \url{https://en.wikipedia.org/wiki/Tanimbili_language}
\item Austronesian Basic Vocabulary Database entry for Tanimbili. \url{https://abvd.eva.mpg.de/austronesian/language.php?id=354}
\item TLOPO (CLLD) language page for Tanibili. \url{https://tlopo.clld.org/languages/tanibili}
\end{itemize}

\section{Religious and Community Resources}
None were found.

\section{Digital and Audio Resources}
\begin{itemize}
\item ANU Open Research repository item, ``Recordings from Solomon Islands, French Polynesia, Papua New Guinea and Australia''. The search result shows only the title, so I could not confirm that it contains Tanibili material. \url{https://openresearch-repository.anu.edu.au/items/4c417461-70d9-4f27-af77-78494a45ee80}
\end{itemize}

\section{Curriculum Development Context}
With about 15 speakers, mostly elderly, this is a revitalisation and documentation case. Practical work would start with the word lists and any archived recordings, in cooperation with speakers and Temotu Province authorities.
\end{document}